% GENERATED, edit the French unit or the translation table

\mxunit{stmg-c08}{Revision and consolidation}
\section{Revision strategy}
\begin{keypoint}
Effective revision alternates retrieval from memory, explicit correction of mistakes and reuse in a problem. Passively rereading the course is not enough: you must attempt a question without a model, compare your answer, then redo the question after a delay.
\end{keypoint}

For each chapter, use the following code: \textbf{V} (mastered), \textbf{O} (shaky), \textbf{R} (to redo). After the diagnosis, work first on an orange or red skill, then check immediately with a similar question.

\begin{center}
\begin{tabularx}{.9\textwidth}{lcccX}
\toprule
Skill & V & O & R & Remedial action \\
\midrule
Choosing a model & $\square$ & $\square$ & $\square$ & Redo a sorting exercise: arithmetic/geometric/exponential.\\
Reading and calculating & $\square$ & $\square$ & $\square$ & Redo the targeted fluency exercises.\\
Writing up and interpreting & $\square$ & $\square$ & $\square$ & Add data, formula, calculation, unit and concluding sentence.\\
\bottomrule
\end{tabularx}
\end{center}

\mxsession{52}{\mxdate*{2027-05-11}}{B: 2 h -- C: 1 h}{Diagnosis of what has been learned}{Revision and consolidation}
\begin{objectives}\begin{itemize}\item Identify what has been learned.\item Classify the mistakes.\item Plan the remediation.\end{itemize}\end{objectives}
\begin{fluency}[6 to 8 minutes, no calculator]\begin{enumerate}\item Write the multiplier for an increase of $\pc{25}$. \mxblank{$1+\dfrac{25}{100}=\num{1,25}$}\item Give the formula for the general term of an arithmetic sequence starting from $u_0$. \mxblank{$u_n=u_0+nr$}\item Solve $10^x=1000$. \mxblank{$x=3$}\item Differentiate $x+\dfrac4x$. \mxblank{$1-\dfrac4{x^2}$}\item State the law of total probability with two branches. \mxblank{$P(B)=P(A)P_A(B)+P(\overline A)P_{\overline A}(B)$}\end{enumerate}\end{fluency}
\begin{prerequisites}Reread the last written record and be able to carry out the basic calculations needed today. Before the investigation, have students rephrase the key words: identify what has been learned and classify the mistakes.\end{prerequisites}

\begin{activity}
Answer the following \mxcount{bilan} questions without notes, two per chapter. Record each error in a grid: knowledge, choice of method, calculation, reading, written presentation. Then choose two personal priorities.
\begin{enumerate}
\item Sequences. The sequence $(u_n)$ is arithmetic with first term $u_0=7$ and common difference $-3$. Calculate $u_{10}$. \mxblank{$-23$}
\item Sequences. The sequence $(v_n)$ is geometric with first term $v_0=200$ and common ratio $\num{1,05}$. Calculate $v_2$. \mxblank{$\num{220,5}$}
\item Exponential functions. State the direction of variation of $f(x)=3\times\num{0,8}^x$ on $\intfo{0}{+\infty}$. \mxblank{decreasing}
\item Exponential functions. A quantity of 500 increases by $\pc{4}$ per year. Write the model $f(x)$ after $x$ years. \mxblank{$500\times\num{1,04}^x$}
\item Common logarithm. Solve $10^x=\num{0,01}$. \mxblank{$x=-2$}
\item Common logarithm. Calculate $\log(1000)-\log(10)$. \mxblank{$2$}
\item Reciprocal function. Solve $\dfrac1x=4$ for $x\neq0$. \mxblank{$\num{0,25}$}
\item Reciprocal function. State the direction of variation of $x\mapsto\dfrac1x$ on $\intoo{0}{+\infty}$. \mxblank{decreasing}
\item Two variables. With the fit $y=2x+3$, estimate $y$ for $x=10$. \mxblank{$23$}
\item Two variables. Give the mean point of the series $x$: 1, 2, 3 and $y$: 4, 6, 8. \mxblank{$G\pt{2;6}$}
\item Conditional probability. $P(A)=\num{0,4}$ and $P_A(B)=\num{0,5}$. Calculate $P(A\cap B)$. \mxblank{$\num{0,2}$}
\item Conditional probability. $P(A\cap B)=\num{0,12}$ and $P(A)=\num{0,3}$. Calculate $P_A(B)$. \mxblank{$\num{0,4}$}
\item Binomial distribution. $X$ follows the distribution $\mathcal B\pt{3;\num{0,5}}$. Calculate $P(X=3)$. \mxblank{$\num{0,125}$}
\item Random variable. $X$ takes the values 0, 1, 2 with probabilities $\num{0,5}$, $\num{0,3}$, $\num{0,2}$. Calculate $E(X)$. \mxblank{$\num{0,7}$}
\mxcounted{bilan}
\end{enumerate}

\end{activity}
\begin{solution}The correction expects a profile of mistakes and a plan for catching up, not just an overall mark. Have students say aloud which piece of data triggers the choice of method.\end{solution}

\begin{summary}The diagnosis is used to choose the remediation: recognise the mistake, recall the useful rule, redo a similar question, then check without help.\end{summary}

\begin{example}
Take the data from the investigation and write the solution in order: definition of the objects, property used, exact or approximate calculation, then a concluding sentence. The model answer is in the teacher edition.
\end{example}
\begin{solution}The correction expects a profile of mistakes and a plan for catching up, not just an overall mark.\end{solution}

\begin{exercise}[essential]
\begin{enumerate}\item Choose between an arithmetic sequence, a geometric sequence and an exponential function in three situations.\item Solve an exponential equation using a logarithm.\item Calculate a probability from a tree, then interpret it.\end{enumerate}
\end{exercise}
\begin{solution}Justify each choice by absolute or relative change; check positivity before taking the logarithm; multiply along a path and conclude with a sentence.\end{solution}

\mxexuse{52}

\begin{homework}Redo practice exercise 3 with the following data: replace the initial value by $380$ and the rate or common ratio by $6$. Write it up in four lines: data, formula, calculation, conclusion.\end{homework}
\begin{solution}The solution depends on the exercise redone. Check above all the consistency of the model, the number of iterations, the rounding and the concluding sentence.\end{solution}

\begin{exitticket}[3 minutes]On a half-sheet: write the main rule of the session, work through a very short example and give a confidence level from 1 to 4. Essential question(s) of the day: 1 and 2.\end{exitticket}

\ifmxteacher
\begin{method}
\textbf{1 h session:} 7 min of fluency; 10 min of pooling and written record; 28 min of essential exercises; 5 min of correction; 3 min of exit ticket.\\
\textbf{2 h session:} add 20 min of investigation, 20 min of formalisation, a second independent exercise and 10 min of comparative correction.\\
\textbf{Differentiation:} give struggling students a four-box sheet (data, model, calculation, conclusion). Ask quick students to produce a counterexample, a justification or a second method.
\end{method}
\fi
\mxsession{53}{\mxdate*{2027-05-13}}{B: 1 h -- C: 2 h}{Reactivation: exponential functions}{Revision and consolidation}
\begin{objectives}\begin{itemize}\item Reactivate exponential models.\item Solve and interpret.\item Check the units.\end{itemize}\end{objectives}
\begin{fluency}[6 to 8 minutes, no calculator]\begin{enumerate}\item Simplify $3x+7-x+2$. \mxblank{$2x+9$}\item Solve $10^x=1000$. \mxblank{$x=3$}\item Differentiate $x+\dfrac4x$. \mxblank{$1-\dfrac4{x^2}$}\item State the law of total probability with two branches. \mxblank{$P(B)=P(A)P_A(B)+P(\overline A)P_{\overline A}(B)$}\item Give the formula for the general term of an arithmetic sequence starting from $u_0$. \mxblank{$u_n=u_0+nr$}\end{enumerate}\end{fluency}
\begin{prerequisites}Reread the last written record and be able to carry out the basic calculations needed today. Before the investigation, have students rephrase the key words: reactivate exponential models and solve and interpret.\end{prerequisites}

\begin{activity}
A value of $\num{18000}$ increases by $\pc{3,5}$ per year. Write the model, calculate the value after eight years and determine when it will exceed $\num{25000}$.
\begin{center}\begin{tikzpicture}\begin{axis}[width=.68\textwidth,height=5cm,axis lines=middle,xmin=-3,xmax=4,ymin=0,ymax=5,grid=major]\addplot[mxteacher,very thick,samples=100,domain=-3:4]{1.4^x};\addplot[mxwarn,very thick,samples=100,domain=-3:4]{.7^x};\end{axis}\end{tikzpicture}\end{center}
\end{activity}
\begin{solution}$V(t)=\num{18000}\times\num{1,035}^t$, $V(8)\approx\num{23703}$. The threshold is exceeded from the tenth year onwards. Have students say aloud which piece of data triggers the choice of method.\end{solution}

\begin{summary}An exponential function models a constant relative change. First identify the initial value and the multiplier.\end{summary}

\begin{example}
Take the data from the investigation and write the solution in order: definition of the objects, property used, exact or approximate calculation, then a concluding sentence. The model answer is in the teacher edition.
\end{example}
\begin{solution}$V(t)=\num{18000}\times\num{1,035}^t$, $V(8)\approx\num{23703}$. The threshold is exceeded from the tenth year onwards.\end{solution}

\begin{exercise}[essential]
\begin{enumerate}\item Choose between an arithmetic sequence, a geometric sequence and an exponential function in three situations.\item Solve an exponential equation using a logarithm.\item Calculate a probability from a tree, then interpret it.\end{enumerate}
\end{exercise}
\begin{solution}Justify each choice by absolute or relative change; check positivity before taking the logarithm; multiply along a path and conclude with a sentence.\end{solution}

\mxexuse{53}

\begin{homework}Redo practice exercise 2 with the following data: replace the initial value by $385$ and the rate or common ratio by $7$. Write it up in four lines: data, formula, calculation, conclusion.\end{homework}
\begin{solution}The solution depends on the exercise redone. Check above all the consistency of the model, the number of iterations, the rounding and the concluding sentence.\end{solution}

\begin{exitticket}[3 minutes]On a half-sheet: write the main rule of the session, work through a very short example and give a confidence level from 1 to 4. Essential question(s) of the day: 1, 2 and 3.\end{exitticket}

\ifmxteacher
\begin{method}
\textbf{1 h session:} 7 min of fluency; 10 min of pooling and written record; 28 min of essential exercises; 5 min of correction; 3 min of exit ticket.\\
\textbf{2 h session:} add 20 min of investigation, 20 min of formalisation, a second independent exercise and 10 min of comparative correction.\\
\textbf{Differentiation:} give struggling students a four-box sheet (data, model, calculation, conclusion). Ask quick students to produce a counterexample, a justification or a second method.
\end{method}
\fi
\mxsession{54}{\mxdate*{2027-05-18}}{B: 2 h -- C: 1 h}{Reactivation: common logarithm (base 10)}{Revision and consolidation}
\begin{objectives}\begin{itemize}\item Reactivate logarithms.\item Solve an equation.\item Check the conditions.\end{itemize}\end{objectives}
\begin{fluency}[6 to 8 minutes, no calculator]\begin{enumerate}\item Calculate $\pc{5}$ of $160$. \mxblank{$8$}\item Differentiate $x+\dfrac4x$. \mxblank{$1-\dfrac4{x^2}$}\item State the law of total probability with two branches. \mxblank{$P(B)=P(A)P_A(B)+P(\overline A)P_{\overline A}(B)$}\item Give the formula for the general term of an arithmetic sequence starting from $u_0$. \mxblank{$u_n=u_0+nr$}\item Solve $10^x=1000$. \mxblank{$x=3$}\end{enumerate}\end{fluency}
\begin{prerequisites}Reread the last written record and be able to carry out the basic calculations needed today. Before the investigation, have students rephrase the key words: reactivate logarithms and solve an equation.\end{prerequisites}

\begin{activity}
Solve $10^x=340$, $\num{1,06}^n=2$ and $\num{1,06}^n\geq2$. For each answer, distinguish the real value, the approximate value and the minimum integer index.

\end{activity}
\begin{solution}$x\approx\num{2,5315}$; $n\approx\num{11,90}$; for the integer threshold, $n=12$. Have students say aloud which piece of data triggers the choice of method.\end{solution}

\begin{summary}The logarithm turns an exponential equation into an equation in the exponent. Conditions: base positive and different from $1$, right-hand side positive.\end{summary}

\begin{example}
Take the data from the investigation and write the solution in order: definition of the objects, property used, exact or approximate calculation, then a concluding sentence. The model answer is in the teacher edition.
\end{example}
\begin{solution}$x\approx\num{2,5315}$; $n\approx\num{11,90}$; for the integer threshold, $n=12$.\end{solution}

\begin{exercise}[essential]
\begin{enumerate}\item Choose between an arithmetic sequence, a geometric sequence and an exponential function in three situations.\item Solve an exponential equation using a logarithm.\item Calculate a probability from a tree, then interpret it.\end{enumerate}
\end{exercise}
\begin{solution}Justify each choice by absolute or relative change; check positivity before taking the logarithm; multiply along a path and conclude with a sentence.\end{solution}

\mxexuse{54}

\begin{homework}Redo practice exercise 3 with the following data: replace the initial value by $390$ and the rate or common ratio by $2$. Write it up in four lines: data, formula, calculation, conclusion.\end{homework}
\begin{solution}The solution depends on the exercise redone. Check above all the consistency of the model, the number of iterations, the rounding and the concluding sentence.\end{solution}

\begin{exitticket}[3 minutes]On a half-sheet: write the main rule of the session, work through a very short example and give a confidence level from 1 to 4. Essential question(s) of the day: 1 and 2.\end{exitticket}

\ifmxteacher
\begin{method}
\textbf{1 h session:} 7 min of fluency; 10 min of pooling and written record; 28 min of essential exercises; 5 min of correction; 3 min of exit ticket.\\
\textbf{2 h session:} add 20 min of investigation, 20 min of formalisation, a second independent exercise and 10 min of comparative correction.\\
\textbf{Differentiation:} give struggling students a four-box sheet (data, model, calculation, conclusion). Ask quick students to produce a counterexample, a justification or a second method.
\end{method}
\fi
\mxsession{55}{\mxdate*{2027-05-20}}{B: 1 h -- C: 2 h}{Reactivation: reciprocal function}{Revision and consolidation}
\begin{objectives}\begin{itemize}\item Reactivate the reciprocal function.\item Carry out a study.\item Sketch and interpret.\end{itemize}\end{objectives}
\begin{fluency}[6 to 8 minutes, no calculator]\begin{enumerate}\item Write the multiplier for an increase of $\pc{10}$. \mxblank{$1+\dfrac{10}{100}=\num{1,10}$}\item State the law of total probability with two branches. \mxblank{$P(B)=P(A)P_A(B)+P(\overline A)P_{\overline A}(B)$}\item Give the formula for the general term of an arithmetic sequence starting from $u_0$. \mxblank{$u_n=u_0+nr$}\item Solve $10^x=1000$. \mxblank{$x=3$}\item Differentiate $x+\dfrac4x$. \mxblank{$1-\dfrac4{x^2}$}\end{enumerate}\end{fluency}
\begin{prerequisites}Reread the last written record and be able to carry out the basic calculations needed today. Before the investigation, have students rephrase the key words: reactivate the reciprocal function and carry out a study.\end{prerequisites}

\begin{activity}
Study $f(x)=x+\dfrac{16}{x}$ on $]0,+\infty[$ and interpret its minimum in an average-cost context. Produce the table of variations and a sketch.
\begin{center}\begin{tikzpicture}\begin{axis}[width=.68\textwidth,height=5cm,axis lines=middle,xmin=-4,xmax=4,ymin=-4,ymax=4,grid=major]\addplot[mxteacher,very thick,samples=80,domain=-4:-.15]{1/x};\addplot[mxteacher,very thick,samples=80,domain=.15:4]{1/x};\end{axis}\end{tikzpicture}\end{center}
\end{activity}
\begin{solution}$f'(x)=\dfrac{(x-4)(x+4)}{x^2}$; minimum $f(4)=8$. Have students say aloud which piece of data triggers the choice of method.\end{solution}

\begin{summary}For a function containing $\dfrac1x$, start by excluding $x=0$, then differentiate, factorise and study the sign before the variations.\end{summary}

\begin{example}
Take the data from the investigation and write the solution in order: definition of the objects, property used, exact or approximate calculation, then a concluding sentence. The model answer is in the teacher edition.
\end{example}
\begin{solution}$f'(x)=\dfrac{(x-4)(x+4)}{x^2}$; minimum $f(4)=8$.\end{solution}

\begin{exercise}[essential]
\begin{enumerate}\item Choose between an arithmetic sequence, a geometric sequence and an exponential function in three situations.\item Solve an exponential equation using a logarithm.\item Calculate a probability from a tree, then interpret it.\end{enumerate}
\end{exercise}
\begin{solution}Justify each choice by absolute or relative change; check positivity before taking the logarithm; multiply along a path and conclude with a sentence.\end{solution}

\mxexuse{55}

\begin{homework}Redo practice exercise 2 with the following data: replace the initial value by $395$ and the rate or common ratio by $3$. Write it up in four lines: data, formula, calculation, conclusion.\end{homework}
\begin{solution}The solution depends on the exercise redone. Check above all the consistency of the model, the number of iterations, the rounding and the concluding sentence.\end{solution}

\begin{exitticket}[3 minutes]On a half-sheet: write the main rule of the session, work through a very short example and give a confidence level from 1 to 4. Essential question(s) of the day: 1, 2 and 3.\end{exitticket}

\ifmxteacher
\begin{method}
\textbf{1 h session:} 7 min of fluency; 10 min of pooling and written record; 28 min of essential exercises; 5 min of correction; 3 min of exit ticket.\\
\textbf{2 h session:} add 20 min of investigation, 20 min of formalisation, a second independent exercise and 10 min of comparative correction.\\
\textbf{Differentiation:} give struggling students a four-box sheet (data, model, calculation, conclusion). Ask quick students to produce a counterexample, a justification or a second method.
\end{method}
\fi
